\documentclass[12pt]{article}
\usepackage[utf8]{inputenc}
\usepackage[T1]{fontenc}
\usepackage{hyperref}
\title{Temoq Language Resources for K-12 Curriculum Development}
\author{Compiled by Claude Sonnet 5 (claude-sonnet-5) for Andr\'as Kornai}
\date{}
\begin{document}
\maketitle

\section{Introduction}
Temoq (ISO 639-3: tmo) is a severely endangered Austroasiatic language of the southern Aslian branch, spoken in Pahang, Peninsular Malaysia. The Temoq people number under 100 (2010 figure). Western Temoq communities mostly use Malay and Semelai, and eastern ones mostly use Jakun, though they still know Temoq. Sources disagree on whether the Temoq are Senoi (the government classification) or Proto-Malay.

\section{Primary Educational Resources}
No primer, textbook or school programme for Temoq turned up in search results.

\section{Linguistic and Reference Materials}
Wikipedia has articles on the language, \url{https://en.wikipedia.org/wiki/Temoq_language}, and on the people, \url{https://en.wikipedia.org/wiki/Temoq_people}. A 2023 Journal of Southeast Asian Studies article discusses the ethnonyms Temoq, Semelai, Semaq Beri and Jakun: \url{https://resolve.cambridge.org/core/journals/journal-of-southeast-asian-studies/article/temoq-semelai-semaq-beri-and-jakun-using-orang-asli-ethnonyms-to-reconstruct-orang-asli-ethnohistory/41199662E31C059C6542ED43BCED9C43}. No Temoq dictionary, grammar or orthography was found.

\section{Religious and Community Resources}
Joshua Project has a people-group profile: \url{https://joshuaproject.net/people_groups/15335}. Asia Harvest has a PDF profile, which I did not open: \url{https://people-groups.asiaharvest.org/Malaysia/Temoq.pdf}. No scripture in Temoq was found.

\section{Digital and Audio Resources}
None found.

\section{Curriculum Development Context}
With under 100 people and elderly-only use of the language, this is a language-documentation and revitalisation situation rather than a school-curriculum one. The first steps would be recording the remaining speakers and building a basic wordlist and orthography with the community.
\end{document}